% Response letter for peer review — point-by-point structure.
\documentclass[11pt]{article}
\usepackage[a4paper, margin=2.4cm]{geometry}
\usepackage{xcolor}
\usepackage{enumitem}
\usepackage{booktabs}
\usepackage{ulem}

\definecolor{replied}{HTML}{0E9F6E}
\newcommand{\reviewer}[1]{\vspace{8pt}\noindent{\large\bfseries #1}\par}
\newcommand{\comment}[2]{\vspace{5pt}\noindent\textbf{\textcolor{blue!60!black}{Comment #1:}} \textit{#2}\par}
\newcommand{\response}[1]{\vspace{3pt}\noindent\textbf{\textcolor{replied}{Response:}} #1\par}

\title{\vspace{-1.2cm}\bfseries Response to Reviewers\\[4pt]
\large Manuscript ID: EXAMPLE-2026-0042}
\author{}
\date{}

\begin{document}
\maketitle

Dear Editor and Reviewers,

We thank the reviewers for their careful reading and constructive
comments. We have revised the manuscript accordingly; all changes are
highlighted in the revised version. Below we address each point.

\reviewer{Reviewer \#1}

\comment{1.1}{The motivation in the Introduction is too brief; please
clarify the research gap.}

\response{We expanded Section~1 (page 1--2, highlighted) to explicitly
state the gap and our contributions. Specifically, we added: ``\ldots''
}

\comment{1.2}{Table 2 lacks error bars; are the improvements significant?}

\response{We added standard deviations over 3 random seeds and a
paired significance test ($p < 0.05$). The revised Table~2 now reports
mean~$\pm$~std.}

\reviewer{Reviewer \#2}

\comment{2.1}{Eq.~(3) should define the symbols $q_i$ and $\lambda$.}

\response{Fixed. Symbols are now defined right after Eq.~(3) on page~4.}

\comment{2.2}{Please discuss the limitation of the self-training loop.}

\response{We added a ``Limitations'' paragraph in Section~5
(page~8) discussing noise accumulation and an early-stopping remedy.}

\vspace{10pt}
\noindent We believe the manuscript has been substantially improved and
hope it is now suitable for publication.

\end{document}
